\centerline{\bf \Huge Первая классическая олимпиада}

\problem{\textbf{1}} В семизначном числе некоторые цифры заменили
звёздочками: *6767**. Известно, что число делится на 8 и на 9. Восстановите число. Перечислите ВСЕ варианты!\\


\problem{\textbf{2}} Ньют хочет перевезти девять фантастических тварей
весом 2, 3, 4, 5, 6, 7, 8, 9 и 10 кг в трех чемоданах, по три твари в каждом. Каждый чемодан должен весить меньше 20 кг. Если вес какой-нибудь твари будет делиться на вес другой твари из того же чемодана, они подерутся. Как Ньюту распределить тварей по чемоданам, чтобы никто не подрался?\\


\problem{\textbf{3}} На улице, став в кружок, разговаривают четыре девочки: Аня, Валя, Галя и Нина. Девочка в зеленом платье (не Аня и не Валя) стоит между девочкой в голубом платье и Ниной. Девочка в белом платье стоит между девочкой в розовом платье и Валей. Какое платье на каждой из девочек?\\

\problem{\textbf{4}} Ягами Лайт заменил в примере на сложение трехзначных чисел цифры буквами: одинаковые – одинаковыми, разные – разными.

У него получилось ЧАС + ПАС = АААУ Каким мог быть исходный пример? Найдите все возможные варианты.\\

\problem{\textbf{5}} На IV Летних математических сборах $X$ пятиклассников. Из них 23 увлекаются шахматами, 18 играют в роблокс, 12 посещают футбольную секцию. 10 шахматистов играют в роблокс, 9 футболистов тоже играют в роблокс и 8 шахматистов ходят на футбол. 6 пятиклассников увлекаются шахматами, играют в роблокс и посещают футбольную секцию. Чему равно $X$, если все пятиклассники увлекаются хотя бы одним из перечисленных занятий?\\

\centerline{\textbf{90 минут}}